\section{Introduction}
    This section introduces the Bellman equation and the Bellman optimality equation,
    using the Bellman operator as a main framework. This part was influenced from \cite{littman1996}.
    Note that we always assume that the state space and action space is finite. Although there are algorithms that can be applied to continuous states or actions (\ref{alg:DQN}),
    for simplicity, we assume finiteness. (and actually, no convergence is guaranteed for DQN.)
\section{The State Value Function and Action Value Function}
    Before diving into the Bellman equation, we introduce an important concept called the \textbf{value of a state} and the \textbf{value of an action}.
    \begin{defn}{State Value}{state-value}
        Let the agent follows the policy $\pi$. The \textbf{state value} $v_\pi(s)$ of state $s$ is defined as the expected return:
        \begin{equation}
            v_\pi(s) = \E[G_t|S_t=s] \label{eq:state-value}
        \end{equation}
        where $G_t$ is the discounted return $G_t=\sum_{k=t+1}^{\infty}\gamma^{k-t-1}R_{k}$.
        The function $v_\pi$ itself is called \textbf{state-value function}.
    \end{defn}
    State values are used to evaluate the policy $\pi$. Policies that generate greater state values are better. Then, how can we
    obtain the state values? The answer is to solve the Bellman equation.
    \begin{defn}{Action Value}{action-value}
        Let $\pi$ be a given policy. Then, the \textbf{action value} of a state-action pair is defined as
        \begin{equation}
            q_\pi(s,a) = \E[G_t|S_t=s,A_t=a].
        \end{equation}
    \end{defn}
    \begin{rem}{Properties of Action Value}{ActionValProp}
        Action value $q_\pi(s,a)$ has following properties.
        \begin{itemize}
            \item $v_\pi=\sum_{a \in \mathcal{A}(s)} \pi(a|s)q_\pi(s,a)$
            \item $q_\pi(s,a) = \sum_{r \in \mathcal{R}(s,a)}p(r|s,a)r + \gamma \sum_{s' \in \mathcal{S}}p(s'|s,a)v_\pi(s')$
        \end{itemize}
    \end{rem}
    \begin{proof}
        \begin{itemize}
            \item Use the law of total expectation \ref{eq:cond-LTE}.
            \item Use the result above and \ref{thm:BellmanEq-SV}.
        \end{itemize}
    \end{proof}
    Actually, the ~\ref{rem:ActionValProp} is the \textbf{Bellman equation} of state values.
\section{Generalized Bellman Equation}
    \begin{defn}{Generalized Bellman Equation}{GeneralizedBellmanEq}
        Let $v(s)$ be a state value function. Then, the \textbf{generalized Bellman equation} is defined as following:
        \begin{equation}
            v(s) = \bigotimes_a \Big(r(s, a) + \gamma \bigoplus_{s'} v(s') \Big) 
        \end{equation}
        where $\bigotimes$ defines how an optimal agent should choose actions, and $\bigoplus$ defines how the value of the next state should be used in assigning value to the current state.

        From here, we can define the \textbf{generalized Bellman operator} $T:\mathbb{R}^\mathcal{S} \to \mathbb{R}^\mathcal{S}$ as following:
        \begin{equation}
            [Tv](s) = \bigotimes_a \Big( r(s,a) + \gamma \bigoplus_{s'}v(s') \Big)
        \end{equation}

        The \textbf{myopic policy} $\pi$ with respect to a value function $v$ is any (stochastic) policy such that $T^\pi v = T v$ where
        \begin{equation}
            [T^\pi v](s) = \sum_a \pi(s|a) \Big( r(s, a) + \gamma \bigoplus_{s'}v(s') \Big)
        \end{equation}
    \end{defn}
    For a myopic policy to exist, we require that the operator $\bigotimes$ satisfy the following property: for all functions $f$ and states $s$,
    $\min_a f(s,a) \leq \bigotimes_a f(s,a) \leq \max_a f(s,a)$ (\cite{littman1996}).
    And we will only use the operator $\Big[\bigotimes_{s'} v\Big](s) = \sum_{s' \in \mathcal{S}} p(s'|s,a) v(s')$, since we only cope with
    MDPs.
\section{Bellman Equation}
    The value function with respect to a policy $\pi$ satisfies the following Bellman equation.
    The Bellman equation is simply a set of linear equations that describe the relationships of each state values.
    Bellman equation can be obtained by using the recursive property of return:
    \begin{align}
        G_t &= R_{t+1} + \gamma R_{t+2} + \gamma^2 R_{t+3} + \gamma^3 R_{t+4} + \cdots \\
            &= R_{t+1} + \gamma (R_{t+2} + \gamma R_{t+3} + \gamma^2 R_{t+4}) + \cdots \\
            &= R_{t+1} + \gamma G_{t+1}.
    \end{align}
    Substituting equation \ref{eq:state-value} with the recursive realationship, we obtain
    \begin{align}
        v_\pi(s) &= \E[G_t|S_t=s] \\
                    &= \E[R_{t+1} + \gamma G_{t+1}|S_t=s] \\
                    &= \E[R_{t+1}|S_t=s] + \gamma \E[G_{t+1}|S_t=s]
    \end{align} 
    Using the law of total expectation \ref{eq:cond-LTE}, we get
    \begin{align}
        \E[R_{t+1}|S_t=s] &= \sum_{a \in \mathcal{A}(s)}\pi(a|s)\E[R_{t+1}|S_t=s,A_t=a] \\
                            &= \sum_{a \in \mathcal{A}(s)}\pi(a|s)\sum_{r \in \mathcal{R}(s,a)}p(r|s,a)r
    \end{align}
    and
    \begin{align}
        \E[G_{t+1}|S_t=s] &= \sum_{s' \in \mathcal{S}}\E[G_{t+1}|S_{t+1}=s',S_t=s]p(s'|s) \\
                            &= \sum_{s' \in \mathcal{S}}\E[G_{t+1}|S_{t+1}=s']p(s'|s) \quad \text{(Markov property)} \\
                            &= \sum_{s' \in \mathcal{S}}v_\pi(s')p(s'|s) \\
                            &= \sum_{s' \in \mathcal{S}}v_\pi(s')\sum_{a \in \mathcal{A}(s)}p(s'|s,a).
    \end{align}
    The final result is given in the following box:
    \begin{thm}{Bellman Equation for State-Value Function}{BellmanEq-SV}
        Let $\pi$ be a policy, and let $v_\pi$ be a state-value function. Then, the following Bellman equation is satisfied:
        \begin{equation}
            v_\pi(s) = \sum_{a \in \mathcal{A}(s)}\pi(a|s) \Big(r(s,a) + \gamma \sum_{s' \in \mathcal{S}}p(s'|s,a)v_\pi(s')\Big)
        \end{equation}
    \end{thm}

    And one can see the operator $\Big[\bigotimes_{s'} v\Big](s) = \sum_{s' \in \mathcal{S}} p(s'|s,a) v(s')$.

    \subsection{Matrix-Vector form of Bellman Equation}
        We can reformulate Bellman equation in terms of matrix-vector form. For this, we rewrite the Bellman equation given in Theorem \ref{thm:BellmanEq-SV} as
        \begin{equation}
            v_\pi(s) = r_\pi(s) + \gamma\sum_{s' \in \mathcal{S}}p_\pi(s'|s)v_\pi(s'),
        \end{equation}
        where
        \begin{align}
            r_\pi(s) &= \sum_{a \in \mathcal{A}(s)}\pi(a|s)\sum_{r \in \mathcal{R}(s,a)}p(r|s,a)r \\
            p_\pi(s'|s) &= \sum_{a \in \mathcal{A}(s)}\pi(a|s)p(s'|s,a).
        \end{align}
        The $r_\pi$ and $p_\pi$ denotes the mean of immediate rewards and probability of transitioning from $s$ to $s'$ under policy $\pi$, respectively.
        Suppose that we indexed all states as $s_i$ with $i = 1, \dots, n$, where $n=|\mathcal{S}|$.
        Let
        \begin{equation}
            v_\pi = \begin{bmatrix}
                    v_\pi(s_1) \\
                    \vdots \\
                    v_\pi(s_n)
                    \end{bmatrix},
            r_\pi = \begin{bmatrix}
                    r_\pi(s_1) \\
                    \vdots \\
                    r_\pi(s_n)
                    \end{bmatrix},
            P_\pi = \begin{bmatrix}
                    p_\pi(s_1|s_1) & p_\pi(s_2|s_1) & \dots & p_\pi(s_n|s_1) \\
                    p_\pi(s_1|s_2) & p_\pi(s_2|s_2) & \dots & p_\pi(s_n|s_2) \\
                    \vdots & \vdots &  & \vdots \\
                    p_\pi(s_1|s_n) & p_\pi(s_2|s_n) & \dots & p_\pi(s_n|s_n)
                    \end{bmatrix}. \label{eq:VecForm}
        \end{equation}
        Then, Theorem \ref{thm:BellmanEq-SV} can be written in the following form:
        \begin{thm}{Matrix-Vector form of Bellman Equation}{MVFormBellmanEq}
        Let $v_\pi$, $r_\pi$ and $P_\pi$ be given as \ref{eq:VecForm}. Then, following is satisfied:
            \begin{equation}
                v_\pi = r_\pi + \gamma P_\pi v_\pi
            \end{equation}
        \end{thm}
        \begin{rem}{Properties of Matrix $P_\pi$}{ProbMatProp}
            Note that the matrix $P_\pi$ given in equation \ref{eq:VecForm} satisfies following properties:
            \begin{itemize}
                \item All the elements are nonnegative;
                \item The sum of the values of every row equals 1. Matrix with this property is called a \emph{stochastic matrix}.
            \end{itemize}
        \end{rem}

\section{Bellman Optimality Equation}
    \subsection{Optimal Policy}
        The reinforcement learning algorithms are all designed to find an optimal policy for
        given enviroment. Then, what is an 'optimal' policy? The definition is given below.
        \begin{defn}{Optimal Policy and Optimal State Value}{optPol}
            A policy $\pi^*$ is optimal if and only if $v_\pi^*(s) \geq v_\pi(s)$ for all $s \in \mathcal{S}$
            and for any other policy $\pi$. The state values of $\pi^*$ are the optimal state values.
        \end{defn}
        \begin{defn}{Better Policy}{betterPol}
            A policy $\pi_1$ is \textbf{better} than a policy $\pi_2$ if and only if
            $v_{\pi_1}(s) \geq v_{\pi_2}(s)$ for all $s \in \mathcal{S}$.
        \end{defn}
        The questions that naturally arise are:
        \begin{itemize}
            \item Existence: Does the optimal policy exists?
            \item Uniqueness: Is the optimal policy unique?
            \item Stochasticity: Is the optimal policy stochastic or deterministic?
            \item Algorithm: How do we obtain the optimal policy and the optimal state values?
        \end{itemize}
    \subsection{Bellman Optimality Equation}
        The optimal policy and optimal state values can be analyised via the \textbf{Bellman optimality equation} (BOE).
        The optimal state values satisfies following Bellman optimality equation. This claim in fact has two logically independent statements:
        \begin{enumerate}
            \item The BOE, viewed purely as a fixed-point equation, has a unique solution (Section \ref{sec:SolveBellmanOptEq});
            \item The unique solution of BOE concides with the optimal stat-value function $v^*$, and the derived greedy policy is optimal (Section \ref{sec:OptPolOptStateValBOE})
        \end{enumerate}
        Note that the Section \ref{sec:SolveBellmanOptEq} does not assume that the state-value function is optimal.
        \begin{defn}{Bellman Optimality Equation}{BellOptEq}
            Let $\Pi$ be a set of possible policies and let $\gamma \in (0,1)$ be a discount rate. Then, the Bellman Optimality Equation is defined as
            \begin{align}
                v(s) &= \max_{a \in \mathcal{A}(s)}\left( r(s,a) + \gamma \sum_{s' \in \mathcal{S}}p(s'|s,a)v(s') \right) \\
                    &= \max_{a \in \mathcal{A}(s)} q(s,a)
            \end{align}
            where $v(s),v(s')$ are unkown variables to be solved.

            The Bellman optimality operator $T^*:\mathbb{R}^\mathcal{S} \to \mathbb{R}^\mathcal{S}$, uses $\bigotimes_a = \max_a$ as its action selection strategy.
            Specifically, the Bellman optimality operator is defined as following:
            \begin{equation}
                [T^* v](s) = \max_{a \in \mathcal{A}(s)}\left( r(s, a) + \gamma \sum_{s' \in \mathcal{S}}p(s'|s,a)v(s') \right)
            \end{equation}
        \end{defn}
        
\section{Solution of The Bellman Equations}
    \subsection{Contraction Mapping Theorem}
        We propose a \textbf{contraction mapping theorem}, the main tool for solving the Bellman equations.
        The contraction mapping theorem is a powerful tool for analyzing general nonlinear equations.
        \begin{defn}{Contraction Mapping}{ContractMap}
            Let $f : \mathbb{R}^d \to \mathbb{R}^d$. Then, the function $f$ is a \textbf{contraction mapping}
            if there exists $\gamma \in (0,1)$ such that
            \begin{equation}
                \|f(x_1) - f(x_2)\| \leq \gamma \|x_1-x_2\|
            \end{equation}
            for all $x_1,x_2 \in \mathbb{R}^d$.
        \end{defn}
        Now we present the contraction mapping theorem.
        \begin{thm}{Contraction Mapping Theorem}{ContractMap}
            Let $f:\mathbb{R}^n \to \mathbb{R}^n$ be a contraction mapping. Then, the following properties hold for equation $x=f(x)$:
            \begin{itemize}
                \item Existence: There exists a fixed point $x^*$ satisfying $f(x^*)=x^*$;
                \item Uniqueness: The fixed point $x^*$ is unique;
                \item Algorithm: The following sequence $\{x_k\}$ converges to the fixed point $x^*$ exponentially, for any initial $x_0$.
                    \begin{equation}
                        x_{k+1} = f(x_k)
                    \end{equation}
            \end{itemize}
        \end{thm}
        \begin{proof}
            \textit{Part 1.} (Convergence) The sequence $\{x_k\}$ defined as $x_{k+1}=f(x_k)$ converges.

            We show that the sequence $\{x_k\}$ is a Cauchy sequence; i.e., there always exists an $N \in \mathcal{N}$ such that $\|x_m-x_n\| < \epsilon$ for any $\epsilon > 0$.
            The Cauchy sequence has a nice property; the Cauchy sequence always converges. This can be proved by using the properties of $\mathcal{R}$, that the $\mathcal{R}$ is a complete set.
            We won't show this here. The following holds:
            \begin{align}
                \|x_{k+1}-x_k\| = &\|f(x_k) - f(x_{k-1})\| \\
                                  &\leq \gamma \|x_k - x_{k-1}\| \\
                                  &\leq \gamma^2 \|x_{k-1}-x_{k-2}\| \\
                                  &\leq \gamma^3 \|x_{k-2}-x_{k-3}\| \\
                                  &\vdots \\
                                  &\leq \gamma^k \|x_1 - x_0\|
            \end{align}
            Using this, we can show that
            \begin{align}
                \|x_m - x_n\| &\leq \|x_m - x_{m-1}\| + \|x_{m-1} - x_{m-2}\| + \cdots + \|x_{n+1} - x_n\| \\
                              &\leq \gamma^{m-1}\|x_1-x_0\| + \gamma^{m-2}\|x_1-x_0\| + \dots + \gamma^{n}\|x_1-x_0\| \\
                              &= \|x_1-x_0\| (\gamma^{m-1} + \gamma^{m-2} + \dots + \gamma^n) \\
                              &= \|x_1-x_0\| \frac{\gamma^n(1-\gamma^{m-n})}{1-\gamma} \\
                              &\leq \|x_1-x_0\| \frac{\gamma^n}{1-\gamma}
            \end{align}
            where $m \geq n$, without loss of generality. Since $\gamma \in (0,1)$, the last line converges to 0. That is, the sequence $\{x_k\}$ is a Cauchy sequence.

            \textit{Part 2.} (Existence) $\lim_{k \to \infty} x_k$ is a solution of fixed-point equation $x=f(x)$.

            Let $x^* = \lim_{k \to \infty} x_k$. Then, following holds.
            \begin{align}
                \|x^*-f(x^*)\| &\leq \|x^*-x_k\| + \|f(x_k)-f(x^*)\| + \|x_k-f(x_k)\| \\
                               &\leq \|x^*-x_k\| + \gamma \|x_k - x^*\| + \|x_k - f(x_k)\|
            \end{align}
            Using the results of \textit{Part 1} and squeeze theorem, we can show that $x^*=f(x^*)$.

            \textit{Part 3.} (Uniqueness) The fixed-point $x^*$ is unique.

            Suppose there is another fixed point $x'$. Then, $\|x^*-x'\| = \|f(x^*)-f(x')\| \leq \gamma \|x^*-x'\|$. Since $\gamma \in (0,1)$, the only possibility is $x^*=x'$.


        \end{proof}

    \subsection{Existence and Uniqueness of the Solution of Bellman Equations}\label{sec:SolveBellmanOptEq}
        Using the contraction mapping theorem, we can prove the following theorem.

        \begin{thm}{Existence and Uniqueness of the Solution of Bellman Equations}{exis-uniq-gen-bellman}
            Let $T:\mathbb{R}^\mathcal{S} \to \mathbb{R}^\mathcal{S}$ be the generalized Bellman operator: $[Tv](s) = \bigotimes_a \Big( r(s,a) + \gamma \bigoplus_{s'}v(s') \Big)$.
            Then, if both $\bigotimes$ and $\bigoplus$ is non-expanding, i.e.,
            $\left|\bigotimes_a f_1(s,a) - \bigotimes_a f_2(s,a)\right| \leq \max_a |f_1(s,a) - f_2(s,a)|$, and $0 \leq \gamma < 1$, the associated Bellman equation
            \begin{equation}
                v(s) = \bigotimes_a \Big(r(s, a) + \gamma \bigoplus_{s'} v(s') \Big) 
            \end{equation} 
            always has solution and the solution is unique.
        \end{thm}
        \begin{proof}
            We show that the operator $T$ is a contraction operator, and use Theorem~\ref{thm:ContractMap}.
            Let $h_v(s,a) = r(s,a) + \gamma \bigoplus_{s'} v(s')$. Then,
            \begin{equation}
                h_{v_1}(s,a) - h_{v_2}(s,a) = \gamma \bigoplus_{s'} \left( v_1(s') - v_2(s') \right)
            \end{equation}
            and using the non-expansion assumption of $\bigoplus$,
            \begin{equation}
                |h_{v_1}(s,a) - h_{v_2}(s,a)| \leq \gamma \max_{s'}|v_1(s') - v_2(s')| = \gamma \|v_1 - v_2\|_\infty
            \end{equation}
            for all $(s, a)$.
            Now, using this and non-expansion assumption of $\bigotimes$, we can get
            \begin{equation}
                |[Tv_1](s) - [Tv_2](s)| = |\bigotimes_a h_{v_1}(s,a) - \bigotimes_a h_{v_2}(s, a)| \leq \max_a |h_{v_1}(s,a) - h_{v_2}(s,a)| \leq \gamma \|v_1 - v_2\|_\infty
            \end{equation}
            for all $s$. Therefore,
            \begin{equation}
                \|[Tv_1](s) - [Tv_2](s)\|_\infty \leq \gamma \|v_1 - v_2\|_\infty
            \end{equation}.
        \end{proof}

        Now, if we prove that $max_a$, $\sum_{s'} \pi(s,a)$ and $\sum_{s' \in \mathcal{S}} p(s'|s,a) v(s')$ is a non-expansion operator,
        the existence and uniqueness of the solution of Bellman equation and Bellman optimality equation is guaranteed, repectively.
        TODO: Show it!

\section{Optimal Policy, Optimal State Value and Bellman Optimality Equation}\label{sec:OptPolOptStateValBOE}
    Now, we should prove that the solution of the Bellman optimality equation is really an optimal state-value function,
    and prove that the greedy policy derived from it is the optimal policy.
    \begin{thm}{Optimality of $v^*$ and $\pi^*$}{OptPolOptStateVal}
        Let $v^*$ be the solution of BOE, and let $\pi^* = \arg\max_{\pi \in \Pi}(r_\pi + \gamma P_\pi v^*)$.
        The, the $v^*$ is an optimal state vaule function, and the $\pi^*$ is an optimal policy. That is, for any policy $\pi$,
        following holds;
        \begin{equation}
            v^* = v_{\pi^*} \geq v_\pi
        \end{equation}
        where $v_\pi$ is the state value function of $\pi$, and $\geq$ is applied elementwise.
    \end{thm}
    \begin{proof}
        For any policy $\pi$, $v_\pi = r_\pi + \gamma P_\pi v_\pi$ holds (Bellman equation). Then,
        \begin{align}
            v^* &= \max_{\tau \in \Pi}(r_\tau + \gamma P_\tau v^*) \\
                &= r_{\pi^*} + \gamma P_{\pi^*} v^* \\
                &\geq r_\pi + \gamma P_\pi v^*
        \end{align}
        and we have
        \begin{equation}
            v^*-v_\pi \geq \gamma P_\pi(v^*-v_\pi).
        \end{equation}
        Using this inequality repeatedly, we have $v^*-v_\pi \geq \gamma^n P_\pi^n (v^*-v_\pi)$ for any $n \in \mathbb{N}$.
        Since $\gamma \in (0,1)$ and $P_\pi$ is a nonnegative matrix with all its elements less than or equal to 1,
        $\lim_{n \to \infty} \gamma^n P_\pi^n (v^*-v_\pi)=0$. That is, $v^*-v_\pi \geq \lim_{n \to \infty} \gamma^n P_\pi^n (v^*-v_\pi)=0$.
    \end{proof}
    The following theorem examines the optimal policy more precisely.
    \begin{thm}{Greedy Optimal Policy}{GreedyOptPol}
        For any $s \in \mathcal{S}$, the following deterministic policy is an optimal policy for solving the BOE.
        \begin{equation}
            \pi^*(a|s) = 
            \begin{cases}
                1, &a = \arg\max_a q^*(a,s) \\
                0, &a \neq \arg\max_a q^*(a,s)
            \end{cases}
        \end{equation}
        where $q^*(s,a) = \sum_{r \in \mathcal{R}}p(r|s,a)r + \gamma \sum_{s' \in \mathcal{S}}p(s'|s,a)v^*(s')$.
    \end{thm}
    \begin{proof}
        From Theorem~\ref{thm:SolutionBellmanOptEq} and~\ref{defn:BellOptEq}, we see that
        \begin{align}
            v^*(s) &= \max_{\pi \in \Pi}\sum_{a \in \mathcal{A}}\pi(a|s)q^*(s,a) \\
                        &\leq \left(\max_{\pi \in \Pi}\sum_{a \in \mathcal{A}}\pi(a|s)\right) \max_{a'} q^*(s,a') \label{tmpeq:2} \\
                        &= \max_{a'}q^*(s,a')
        \end{align}
        and the given deterministic policy $\pi^*$ satisfies the equality of inequality~\ref{tmpeq:2}.
    \end{proof}
    The important truth is that the optimal state value is unique but the optimal policy is not. The Theorem~\ref{thm:GreedyOptPol}
    only shows that there always exists at least one deterministic optimal policy. There still can be other optimal policies, too, stochastic or whatever.

\section{Factors that Influence Optimal Policies}
    Now, we analyze the factors of BOE that influences optimal policies. The factors are: reward and discount rate.
    \subsection{Reward}
        Changing reward can lead to slightly different behavior of policies. For example, in the Figure 3.4(a) of \cite{Zhao2025} (gridworld), one can see that
        the agent tries to move to the goal even though there are forbidden areas. The forbidden area gives the negative reward, but the cumulative reward is greater than
        going around the forbidden areas. This is why reward (signal) is important.
    \subsection{Discount Rate}
        High dicount rate leads to far-sighted agent, and low discount rate leads to short-sighted agent. In the extreme, when $\gamma$, the agent only
        take into account the immediate reward.
    \subsection{Invariance of the Optimal Policy}
        Changing the reward may lead to different optimal policies, but this is not always the case. If we apply the same affine transformation
        to all rewards, the optimal policy remains the same.
        \begin{thm}{Optimal Policy Invariance}{OptPolInvariance}
            Let $v^* \in \mathbb{R}^{|\mathcal{S}|}$ be the optimal state value satisfying $v^* = \max_{\pi \in \Pi}{(r_\pi + \gamma P_\pi v^*)}$.
            If every reward $r \in \mathcal{R}$ goes through the same affine transformation $\alpha * r + \beta$, where $\alpha, \beta \in \mathbb{R}$ and $\alpha > 0$,
            the corresponding optimal state value $v'$ is also an affine transformation of $v^*$:
            \begin{equation}
                v' = \alpha v^* + \frac{\beta}{1 - \gamma}\mathbb{1}
            \end{equation}
            where $\gamma \in (0,1)$ is the discount rate and $\mathbb{1} = [1,\dots,1]^\top$.
            Due to this property, the optimal policy stays invariant.
        \end{thm}